\chapter{Supplement: Generalised Linear Models and Model Assessment}

\begin{nuance}{place in the course}
This supplementary chapter extends the regression of Chapters~5--6 in two directions the textbook does not cover.
\begin{itemize}
  \item Responses that are not normal and continuous (yes/no outcomes, counts) are handled by \emph{generalised linear models}. Logistic regression is the most important.
  \item We ask how well a fitted model will predict \emph{new} data. This leads to cross-validation, information criteria and regularisation (lasso and ridge).
\end{itemize}
\end{nuance}

\section{Why ordinary regression is not enough}

Suppose the response is binary, $Y\in\{0,1\}$ (a customer converts or not). The linear model $\Exp[Y\mid x]=\beta_0+\beta_1x$ has three problems:
\begin{itemize}
  \item it can predict probabilities outside $[0,1]$;
  \item the errors have non-constant variance ($\Var[Y]=p(1-p)$ depends on the mean);
  \item the errors are far from normal.
\end{itemize}
Counts have the same trouble: negative predictions, and variance that grows with the mean.

The fix: model a \emph{transformation} of the mean as a linear function of the regressors, and use the correct distribution for the response.

\section{Generalised linear models}

\begin{definition}[Generalised linear model]
A generalised linear model (GLM) has three parts:
\begin{enumerate}[label=(\arabic*)]
  \item a \emph{random component}: independent $Y_i$ from an exponential-family distribution (normal, Bernoulli/binomial, Poisson, gamma, \dots) with mean $\mu_i=\Exp[Y_i]$;
  \item a \emph{linear predictor} $\eta_i=\beta_0+\beta_1x_{i1}+\dots+\beta_kx_{ik}=\mathbf{x}_i\T\boldsymbol\beta$;
  \item a \emph{link function} $g$ with $g(\mu_i)=\eta_i$.
\end{enumerate}
\end{definition}

\begin{center}
\begin{tabular}{llll}
\toprule
Response & Distribution & Canonical link $g(\mu)$ & Inverse link\\
\midrule
continuous & normal & identity, $\mu$ & $\eta$\\
binary & Bernoulli & logit, $\ln\frac\mu{1-\mu}$ & $\dfrac1{1+e^{-\eta}}$\\
counts & Poisson & log, $\ln\mu$ & $e^\eta$\\
positive continuous & gamma & inverse, $1/\mu$ & $1/\eta$\\
\bottomrule
\end{tabular}
\end{center}

Ordinary multiple regression (Chapter~6) is the GLM with a normal response and identity link.

\section{Logistic regression}

For $Y_i\sim\mathrm{Bernoulli}(p_i)$, the \emph{logistic regression} model is
\[ \ln\frac{p_i}{1-p_i}=\mathbf{x}_i\T\boldsymbol\beta,\qquad\text{that is,}\qquad p_i=\frac1{1+e^{-\mathbf{x}_i\T\boldsymbol\beta}}. \]
$p/(1-p)$ is the \emph{odds} of success. $\ln\frac p{1-p}$ is the \emph{log-odds} (logit).

\begin{theorem}[Interpretation of the coefficients]\label{thm:logit-interp}
A one-unit increase in $x_j$, with the other regressors fixed, adds $\beta_j$ to the log-odds. It \emph{multiplies the odds by} $e^{\beta_j}$, the \emph{odds ratio}. For a binary regressor, $e^{\beta_j}$ is the odds in the group with $x_j=1$ divided by the odds in the group with $x_j=0$.
\end{theorem}

\begin{example}
In a conversion model, $\hat\beta_1=0.70$ for the indicator ``saw the new design''.
\begin{itemize}
  \item Odds ratio: $e^{0.70}=2.01$. The odds of converting are about twice as high.
  \item If the baseline probability is $10\%$ (odds $1/9$), the new odds are $2/9$, so $p=2/11=18\%$.
  \item The effect on the \emph{probability} depends on the baseline. The effect on the odds is constant.
\end{itemize}
\end{example}

\subsection{Maximum likelihood estimation}

For independent Bernoulli observations, the log-likelihood is $\ell(\boldsymbol\beta)=\sum_i\big[y_i\eta_i-\ln(1+e^{\eta_i})\big]$ with $\eta_i=\mathbf{x}_i\T\boldsymbol\beta$.

\begin{theorem}[Score and Hessian for logistic regression]\label{thm:logit-score}
With $p_i=1/(1+e^{-\eta_i})$,
\[ \nabla\ell(\boldsymbol\beta)=X\T(\mathbf{y}-\mathbf{p}),\qquad\nabla^2\ell(\boldsymbol\beta)=-X\T WX,\qquad W=\diag\big(p_i(1-p_i)\big). \]
$X\T WX$ is positive semidefinite (positive definite when $X$ has full column rank and $0<p_i<1$). So $\ell$ is concave, and the MLE, when it exists, is unique.
\end{theorem}

\begin{proof}
$\partial\ell/\partial\beta_j=\sum_i(y_i-p_i)x_{ij}$, since $\partial\ln(1+e^{\eta})/\partial\eta=p$.

Differentiate again: $\partial p_i/\partial\beta_k=p_i(1-p_i)x_{ik}$. This gives the Hessian $-\sum_ip_i(1-p_i)\mathbf{x}_i\mathbf{x}_i\T=-X\T WX$.

For any $\mathbf{v}$: $\mathbf{v}\T X\T WX\mathbf{v}=\sum_iw_i(\mathbf{x}_i\T\mathbf{v})^2\ge0$.
\end{proof}

The likelihood equations $X\T(\mathbf{y}-\mathbf{p})=\mathbf{0}$ have no closed form. Solve them by \emph{Newton's method}, which here is called \emph{iteratively reweighted least squares} (IRLS):
\[ \boldsymbol\beta^{(t+1)}=\boldsymbol\beta^{(t)}+(X\T W^{(t)}X)^{-1}X\T(\mathbf{y}-\mathbf{p}^{(t)})=(X\T W^{(t)}X)^{-1}X\T W^{(t)}\mathbf{z}^{(t)}. \]
The \emph{working response} is $z_i=\eta_i+(y_i-p_i)/(p_i(1-p_i))$. Each iteration is a weighted least squares fit (Chapter~6), with weights $W^{(t)}$ updated from the current fit.

\begin{pitfall}{perfect separation}
Suppose some combination of the regressors separates the $0$'s from the $1$'s perfectly. Then the likelihood keeps increasing as the coefficients grow. The MLE does not exist: some $\hat\beta_j\to\pm\infty$, and software warns about convergence.

Remedies: penalised (ridge-type) logistic regression, a prior (Bayesian logistic regression), or dropping the offending regressor.
\end{pitfall}

\subsection{Inference}

By Theorem~\ref{thm:mle-asym}, in large samples $\hat{\boldsymbol\beta}\approx\Normal\big(\boldsymbol\beta,(X\T WX)^{-1}\big)$, with $W$ evaluated at $\hat{\boldsymbol\beta}$. So:
\begin{itemize}
  \item \textbf{Wald test} for $H_0:\beta_j=0$: $z=\hat\beta_j/\se(\hat\beta_j)$, with $\se(\hat\beta_j)=\sqrt{[(X\T WX)^{-1}]_{jj}}$, compared with $\Normal(0,1)$. A $95\%$ interval for the odds ratio is $\exp\big(\hat\beta_j\pm1.96\,\se(\hat\beta_j)\big)$.
  \item \textbf{Likelihood ratio test} for nested models: $-2\big(\ell_{\mathrm{reduced}}-\ell_{\mathrm{full}}\big)\approx\chi^2_r$ under $H_0$, where $r$ is the number of coefficients set to zero (Wilks' theorem, Chapter~3). This is the GLM version of the extra-sum-of-squares $F$-test.
\end{itemize}

\begin{definition}[Deviance]
The \emph{deviance} is $D=2\big(\ell_{\mathrm{saturated}}-\ell(\hat{\boldsymbol\beta})\big)$. It is twice the log-likelihood gap to a saturated model that fits each observation exactly. For ungrouped Bernoulli data, $D=-2\ell(\hat{\boldsymbol\beta})$. It generalises the residual sum of squares: the change in deviance between nested models is the likelihood ratio statistic.
\end{definition}

\begin{example}[Likelihood ratio test]
Adding two regressors lowers the deviance from $412.6$ to $405.1$. The statistic is $7.5$. Compare with $\chi^2_{0.05,2}=5.99$. The two regressors are jointly significant at $5\%$ ($P=0.023$).
\end{example}

\subsection{Classification and evaluation}

A fitted logistic model gives probabilities $\hat{p}_i$. A decision rule calls a case positive if $\hat{p}_i>c$. The usual choice is $c=0.5$, but the threshold should reflect the costs of the two kinds of error.

Ways to measure performance:
\begin{itemize}
  \item The confusion matrix: sensitivity (true positive rate) and specificity (true negative rate).
  \item The \emph{ROC curve}: true positive rate against false positive rate as $c$ varies.
  \item The area under it (\emph{AUC}): the probability that a random positive has a higher score than a random negative.
  \item \emph{Calibration}: among cases with $\hat{p}\approx0.3$, are about $30\%$ positive?
\end{itemize}

\section{Poisson regression}

For counts $Y_i\sim\mathrm{Poisson}(\mu_i)$ with $\ln\mu_i=\mathbf{x}_i\T\boldsymbol\beta$:
\begin{itemize}
  \item An \emph{offset} $\ln t_i$ (exposure time or population size) models a \emph{rate}.
  \item The score is $X\T(\mathbf{y}-\boldsymbol\mu)$. The information matrix is $X\T WX$ with $W=\diag(\mu_i)$.
  \item $e^{\beta_j}$ is the factor by which the expected count is multiplied for a unit increase in $x_j$ (a \emph{rate ratio}).
  \item The Poisson model forces $\Var[Y]=\mu$. If the data are \emph{overdispersed} ($\Var[Y]>\mu$, estimated by Pearson's statistic over its degrees of freedom), use a quasi-Poisson or negative binomial model. Otherwise standard errors are too small and tests too liberal.
\end{itemize}

\section{Model assessment: how well will the model predict?}

A model fitted to data fits those data better than it fits new data, because it adapts to the noise.

\begin{definition}[Training error and test error]
\begin{itemize}
  \item The \emph{training error} is the average loss on the data used to fit the model.
  \item The \emph{test (generalisation) error} is the expected loss on new observations from the same population.
\end{itemize}
For squared error and a fitted function $\hat{f}$, $\mathrm{Err}=\Exp[(Y_0-\hat{f}(\mathbf{x}_0))^2]$ for a new point $(\mathbf{x}_0,Y_0)$.
\end{definition}

\begin{theorem}[Bias--variance decomposition]\label{thm:biasvar-pred}
Let $Y_0=f(\mathbf{x}_0)+\varepsilon_0$ with $\Exp[\varepsilon_0]=0$ and $\Var[\varepsilon_0]=\sigma^2$, independent of the training data. Then
\[ \Exp\big[(Y_0-\hat{f}(\mathbf{x}_0))^2\big]=\underbrace{\sigma^2}_{\text{irreducible noise}}+\underbrace{\big(\Exp[\hat{f}(\mathbf{x}_0)]-f(\mathbf{x}_0)\big)^2}_{\text{bias}^2}+\underbrace{\Var[\hat{f}(\mathbf{x}_0)]}_{\text{variance}}. \]
\end{theorem}

\begin{proof}
Add and subtract $f(\mathbf{x}_0)$ and $\Exp[\hat{f}(\mathbf{x}_0)]$, and expand. The cross terms vanish because $\varepsilon_0$ has mean zero and is independent of $\hat{f}$, and $\Exp[\hat{f}-\Exp[\hat{f}]]=0$. This is the decomposition $\MSE=\Var+\bias^2$ of Chapter~1 applied to prediction, plus the unavoidable noise $\sigma^2$.
\end{proof}

\begin{itemize}
  \item More flexible models (more regressors, higher-degree polynomials, weaker regularisation) have lower bias and higher variance.
  \item Simple models have high bias and low variance.
  \item Test error is lowest at an intermediate complexity. This is the bias--variance trade-off.
  \item Training error keeps falling as complexity grows, so it is useless for choosing a model.
\end{itemize}

\begin{theorem}[Optimism of the training error]\label{thm:optimism}
For a least squares fit with $p$ parameters and error variance $\sigma^2$, the expected test error on new responses at the \emph{same} regressor values exceeds the expected training error by $2p\sigma^2/n$ per observation:
\[ \Exp[\mathrm{Err}_{\mathrm{in}}]=\Exp[\mathrm{err}_{\mathrm{train}}]+\frac{2p}n\sigma^2. \]
\end{theorem}

\begin{proof}
For the hat matrix $H$ (Chapter~6), $\hat{\mathbf{y}}=H\mathbf{y}$ and $\sum_i\Cov(\hat{y}_i,y_i)=\sigma^2\tr H=p\sigma^2$. The standard identity $\Exp[\mathrm{Err}_{\mathrm{in}}-\mathrm{err}_{\mathrm{train}}]=\frac2n\sum\Cov(\hat{y}_i,y_i)$ gives the result.
\end{proof}

This motivates Mallows' $C_p=\mathrm{err}_{\mathrm{train}}\cdot n/\hat\sigma^2+2p-n$ (Chapter~6) and the information criteria below.

\begin{definition}[AIC and BIC]
For a model with $p$ parameters and maximised log-likelihood $\hat\ell$ on $n$ observations,
\[ \mathrm{AIC}=-2\hat\ell+2p,\qquad\mathrm{BIC}=-2\hat\ell+p\ln n. \]
Smaller is better.
\begin{itemize}
  \item AIC aims at predictive accuracy (it estimates out-of-sample log-loss).
  \item BIC penalises complexity more for $n>7$. It is consistent for picking the true model among those considered, when one of them is true.
\end{itemize}
\end{definition}

\section{Cross-validation}

Information criteria rely on distributional assumptions. \emph{Cross-validation} (CV) estimates the test error directly, by repeatedly holding out part of the data.

\begin{definition}[$K$-fold cross-validation]
Split the data randomly into $K$ roughly equal folds. For $k=1,\dots,K$: fit the model on the other $K-1$ folds and compute the loss on fold $k$. The CV estimate is the average held-out loss over all $n$ observations. $K=5$ or $10$ is typical. $K=n$ is \emph{leave-one-out} cross-validation (LOOCV).
\end{definition}

\begin{theorem}[LOOCV shortcut for linear regression]\label{thm:loocv}
For least squares with hat-matrix diagonal $h_{ii}$ and residuals $e_i$, the leave-one-out prediction error for observation $i$ is $e_i/(1-h_{ii})$. So
\[ \mathrm{LOOCV}=\frac1n\sum_{i=1}^n\Big(\frac{e_i}{1-h_{ii}}\Big)^2=\frac{PRESS}n, \]
from a single fit.
\end{theorem}

This is the PRESS statistic of Chapter~6. Practical rules:
\begin{itemize}
  \item Do any data-dependent preprocessing, feature selection or tuning \emph{inside} each training fold. Otherwise the CV estimate is optimistic.
  \item For time-ordered data, use forward-chaining splits, not random folds.
  \item For grouped data (several rows per user), split by group.
\end{itemize}

\begin{interview}{training, validation and test sets}
A clean workflow:
\begin{itemize}
  \item Training data: fit models.
  \item Validation set (or CV): choose among models and tune hyperparameters.
  \item A final untouched \emph{test set}, used once: report the performance of the chosen model.
\end{itemize}
Reusing the test set to pick models turns it into a validation set and gives optimistic estimates.
\end{interview}

\section{Regularisation: ridge and lasso}

When $p$ is large relative to $n$, least squares overfits. Penalising the size of the coefficients trades a little bias for a big drop in variance (Theorem~\ref{thm:biasvar-pred}).

\begin{definition}[Ridge and lasso]
For $\lambda\ge0$ (regressors standardised, intercept unpenalised):
\begin{align*}
\hat{\boldsymbol\beta}^{\mathrm{ridge}}&=\arg\min_{\boldsymbol\beta}\ \tfrac12\norm{\mathbf{y}-X\boldsymbol\beta}^2+\tfrac\lambda2\norm{\boldsymbol\beta}_2^2,\\
\hat{\boldsymbol\beta}^{\mathrm{lasso}}&=\arg\min_{\boldsymbol\beta}\ \tfrac12\norm{\mathbf{y}-X\boldsymbol\beta}^2+\lambda\norm{\boldsymbol\beta}_1.
\end{align*}
\end{definition}

\begin{itemize}
  \item Ridge: $(X\T X+\lambda I)^{-1}X\T\mathbf{y}$ (Chapter~6; Linear Algebra, Chapter~17). Coefficients shrink smoothly but are never exactly zero.
  \item Lasso: no closed form in general. For an \emph{orthonormal} design ($X\T X=I$) it is explicit.
\end{itemize}

\begin{theorem}[Lasso with orthonormal design]\label{thm:lasso}
If $X\T X=I$, and $\hat\beta_j^{\mathrm{OLS}}$ are the least squares coefficients, then
\[ \hat\beta_j^{\mathrm{lasso}}=\sgn\big(\hat\beta_j^{\mathrm{OLS}}\big)\big(\abs{\hat\beta_j^{\mathrm{OLS}}}-\lambda\big)_+\quad\text{(soft thresholding)},\qquad\hat\beta_j^{\mathrm{ridge}}=\frac{\hat\beta_j^{\mathrm{OLS}}}{1+\lambda}. \]
\end{theorem}

\begin{proof}
With $X\T X=I$, $\tfrac12\norm{\mathbf{y}-X\boldsymbol\beta}^2=\text{const}+\tfrac12\sum_j(\beta_j-\hat\beta_j^{\mathrm{OLS}})^2$. The problem splits into $p$ scalar problems $\min_\beta\tfrac12(\beta-b)^2+\lambda\abs{\beta}$.

Their solution is $\sgn(b)(\abs{b}-\lambda)_+$. For $\beta>0$ the derivative $\beta-b+\lambda$ vanishes at $b-\lambda$, valid when $b>\lambda$. And $\beta=0$ is optimal exactly when $\abs{b}\le\lambda$.

The ridge case differentiates $\tfrac12(\beta-b)^2+\tfrac\lambda2\beta^2$.
\end{proof}

So:
\begin{itemize}
  \item The lasso sets small coefficients exactly to zero (\emph{variable selection}) and shrinks the others by the constant $\lambda$.
  \item Ridge shrinks every coefficient by the same factor $1/(1+\lambda)$.
  \item Geometrically, the contours of the squared error first touch the diamond $\norm{\boldsymbol\beta}_1\le t$ at a corner, where some coordinates vanish. The round ball $\norm{\boldsymbol\beta}_2\le t$ has no corners (Linear Algebra, Chapter~1).
  \item Choose $\lambda$ by cross-validation. It is the complexity knob of the bias--variance trade-off.
  \item The same penalties work for logistic and other GLMs: add them to the negative log-likelihood.
\end{itemize}

\begin{summarybox}
\begin{itemize}
  \item A GLM is an exponential-family response, a linear predictor $\eta=\mathbf{x}\T\boldsymbol\beta$, and a link $g(\mu)=\eta$. Logistic regression: logit link, $e^{\beta_j}$ is an odds ratio. Poisson regression: log link, $e^{\beta_j}$ is a rate ratio.
  \item Logistic MLE: score $X\T(\mathbf{y}-\mathbf{p})$, Hessian $-X\T WX$ (concave), solved by IRLS. Use Wald and likelihood-ratio tests. Deviance generalises the residual sum of squares. Beware perfect separation and overdispersion.
  \item Test error $=\sigma^2+\text{bias}^2+\text{variance}$. Training error is optimistic by $2p\sigma^2/n$. Compare models by AIC, BIC or cross-validation ($K$-fold; LOOCV $=PRESS/n$ for linear regression).
  \item Ridge shrinks smoothly. Lasso soft-thresholds and selects variables. Tune $\lambda$ by cross-validation, with all preprocessing inside the folds and a final untouched test set.
\end{itemize}
\end{summarybox}

\section{Exercises}

\begin{exercise}
In a logistic regression, $\hat\beta_0=-2.0$ and $\hat\beta_1=0.5$ for a regressor $x$. Compute the estimated probability of success at $x=0,2,4$ and the odds ratio for a unit increase in $x$.
\end{exercise}

\begin{exercise}
Derive the score and the Hessian of the Poisson regression log-likelihood $\ell(\boldsymbol\beta)=\sum_i[y_i\mathbf{x}_i\T\boldsymbol\beta-e^{\mathbf{x}_i\T\boldsymbol\beta}-\ln y_i!]$, and write one IRLS update.
\end{exercise}

\begin{exercise}
A $2\times2$ table gives: exposed, $30$ cases of $100$; unexposed, $15$ cases of $100$. Compute the odds ratio and show that it equals $e^{\hat\beta_1}$ for the logistic regression on the exposure indicator. Give a $95\%$ confidence interval using the standard error $\sqrt{1/a+1/b+1/c+1/d}$ of $\ln\widehat{OR}$.
\end{exercise}

\begin{exercise}
Two models for $n=200$ observations have maximised log-likelihoods $-310.2$ ($p=3$) and $-305.9$ ($p=6$). Compute AIC and BIC for each; which would each criterion choose, and does the likelihood ratio test agree?
\end{exercise}

\begin{exercise}
Prove Theorem~\ref{thm:loocv} for simple linear regression through the origin ($\hat\beta=\sum x_iy_i/\sum x_i^2$) and check it numerically for the data $(1,2),(2,3),(3,5)$.
\end{exercise}

\begin{exercise}
Show that with orthonormal design the lasso solution is zero exactly when $\abs{\hat\beta_j^{\mathrm{OLS}}}\le\lambda$, and compute the lasso and ridge coefficients for $\hat\beta^{\mathrm{OLS}}=(3.0,\ 0.8,\ -1.5,\ 0.2)$ and $\lambda=1$.
\end{exercise}

\begin{exercise}
Explain why the training error always decreases as polynomial degree increases while the cross-validation error eventually increases, and sketch both curves. What polynomial degree would you choose from a CV curve that is flat beyond degree $3$?
\end{exercise}

\begin{exercise}
A classifier scores $5$ positives $0.9,0.8,0.55,0.4,0.3$ and $5$ negatives $0.7,0.5,0.35,0.2,0.1$. Compute the AUC as the probability that a random positive outscores a random negative, and the sensitivity and specificity at threshold $0.5$.
\end{exercise}
